\section{From one-two reductions to general complementary depth}

For positive indices we abbreviate
\begin{equation}\label{complement:eq:mpl}
\Li_{s_1,\ldots,s_d}(z)
 :=\Li_{s_1,\ldots,s_d}(z,1,\ldots,1)
 =\sum_{n_1>\cdots>n_d\ge1}
 \frac{z^{n_1}}{n_1^{s_1}\cdots n_d^{s_d}},\qquad |z|<1.
\end{equation}
The indices decrease in the summation order, as in the manuscript.
Weight is $w=s_1+\cdots+s_d$ and series depth is $d$.
An empty index denotes the multiplicative identity.
We use convergent multiple zeta values
$\zeta(\boldsymbol s)=\Li_{\boldsymbol s}(1)$ when $s_1\ge2$.

Let
\[
 \omega_0(t)=\frac{dt}{t},\qquad
 \omega_1(t)=\frac{dt}{1-t}.
\]
Words are written \emph{outer letter first}. Thus
\begin{equation}\label{complement:eq:word}
H_{0^{s_1-1}1\cdots0^{s_d-1}1}(z)=\Li_{s_1,\ldots,s_d}(z).
\end{equation}
In particular $H_{01}=\Li_2$, not its negative or a reversed word.
The recursion is
$H_{av}(z)=\int_0^z\omega_a(t)H_v(t)$ whenever this is an ordinary
convergent integral. Pure-zero words are defined by
$H_{0^r}(z)=\log^r z/r!$, and trailing-zero words are obtained by shuffle
regularization. Every word ending in $1$ is convergent at its initial
endpoint and has the ordinary interpretation in \eqref{complement:eq:word}.

Throughout the functional identities we use
\begin{equation}\label{complement:eq:domain}
 D=\C\setminus[0,\infty),\qquad
 q=\frac1{1-z},\qquad L=\log q=-\log(1-z),\qquad z\in D,
\end{equation}
with principal logarithms and principal polylogarithms. The image of $D$
avoids both $(-\infty,0]$ and $[1,\infty)$ in the $q$-plane, so the
hyperlogarithms with possible trailing zeros have unambiguous branches.
The origin is a limiting point, not an included point of $D$.

We first prove formulas on $z<0$, where $0<q<1$ and all logarithms in
$q$ are real. Analytic continuation on the connected slit domain $D$
then proves the stated complex identities. This method avoids an
unstated choice of a value on the polylogarithm cut. It includes both
$z=i$ and $z=\rho=e^{2\pi i/3}$:
\begin{align}
 z=i:&\quad q=\frac{1+i}{2},&L&=-\frac12\log2+\frac{i\pi}{4},\label{complement:eq:gausspoint}\\
 z=\rho:&\quad q=\frac12+\frac{i\sqrt3}{6},&L&=-\frac12\log3+\frac{i\pi}{6}.\label{complement:eq:eispoint}
\end{align}
Their transformed moduli are $2^{-1/2}$ and $3^{-1/2}$ respectively.

\begin{definition}[Depth with logarithmic coefficients]
A term $L^j\Li_{\boldsymbol a}(q)\zeta(\boldsymbol b)$ has
\emph{reduced depth} $\dep(\boldsymbol a)+\dep(\boldsymbol b)$.
The factor $L^j$ is a coefficient, not an additional nested sum.
\end{definition}
This is an explicit convention for the following reduction. It is not
an identification with every motivic depth filtration, and does not
turn arbitrary products of deep factors into depth-one objects.

\section{An explicit complementary-depth theorem}

\subsection{Endpoint transport with convergent constants}
For a word $v$ ending in $0$, write $J_v(q)$ for the iterated integral
from $1$ to $q$ using the same forms and the same outer-first convention.
These integrals converge at $1$: the innermost form is $\omega_0$, which
is regular there; preceding $\omega_1$ integrations do not destroy
integrability. Set $J_\emp=1$.

\begin{lemma}[Connection formula]\label{complement:lem:connection}
For every nonempty word $v$ ending in $0$ and $0<q<1$,
\begin{equation}\label{complement:eq:connection}
 J_v(q)=\sum_{v=uv'}(-1)^{|v'|}H_u(q)H_{\operatorname{rev}(v')}(1).
\end{equation}
The sum includes empty prefixes and suffixes. Every nonempty endpoint
word $\operatorname{rev}(v')$ starts in $0$, so its value at $1$ is
convergent after the usual trailing-zero normalization at $0$.
\end{lemma}
\begin{proof}
We give the endpoint justification rather than silently substituting into
a divergent path identity. Fix a maximum word length $N$ and work in the
truncated noncommutative algebra on $e_0,e_1$. Let
$F(q)=\sum_{|u|\le N}H_u(q)u$. It satisfies
\[
 F'(q)=\left(\frac{e_0}{q}+\frac{e_1}{1-q}\right)F(q).
\]
Put $\ell=-\log(1-q)$. The function $e^{-\ell e_1}F(q)$ has a finite
coefficientwise limit $Z$ as $q\to1^-$: its differential equation has
coefficient $e^{-\ell e_1}e_0e^{\ell e_1}/q$, a finite polynomial in
$\ell$ at each truncated weight, and this is integrable at $1$.
Consequently
$F(q)=e^{\ell e_1}Z+O((1-q)\log^M(1-q))$
coefficientwise, for a sufficiently large finite $M$.

Shuffle identities say that $F$ is group-like; the same holds for $Z$.
The inverse coefficient in a group-like word series is
$(Z^{-1})_{v'}=(-1)^{|v'|}Z_{\operatorname{rev}(v')}$, by the shuffle
antipode, or equivalently by reversal of an iterated path.
The coefficients of $R(q)=F(q)Z^{-1}$ solve the same differential
recursion and tend to zero at $1$ for every nonempty word ending in $0$.
Induction on word length therefore identifies them with the convergent
integrals $J_v(q)$.

If $v'$ is a nonempty suffix of $v$, its reverse begins in $0$.
The corresponding coefficient of $Z$ is simply
$H_{\operatorname{rev}(v')}(1)$: multiplication by
$e^{-\ell e_1}$ only changes words beginning with $1$.
Taking the coefficient of $v$ in $F(q)Z^{-1}$ proves
\eqref{complement:eq:connection}. Since $N$ was arbitrary, this proves every finite
word identity.
\end{proof}

\subsection{The word formula and the depth bound}
For a word $w$ ending in $1$, let $K$ be its number of zeros. Replace
each original $1$ by $0$, and replace each original $0$ independently
by either $0$ or $1$. Denote the resulting $2^K$ words by $\tau_A(w)$,
where $A$ records the zero positions replaced by $1$.
Every $\tau_A(w)$ ends in $0$.

\begin{theorem}[Complementary depth]\label{complement:thm:complement}
For $z\in D$ and every nonempty word $w$ ending in $1$,
\begin{equation}\label{complement:eq:master}
 H_w(z)=(-1)^K\sum_A\ \sum_{\tau_A(w)=uv}
 (-1)^{|v|}H_u(q)H_{\operatorname{rev}(v)}(1).
\end{equation}
After trailing-zero normalization, this is a finite rational linear
combination of terms
\begin{equation}\label{complement:eq:termshape}
 L^j\Li_{\boldsymbol a}(q)\zeta(\boldsymbol b),\qquad
 j+\wt(\boldsymbol a)+\wt(\boldsymbol b)=|w|,
 \quad \dep(\boldsymbol a)+\dep(\boldsymbol b)\le K.
\end{equation}
Nonempty zeta indices are convergent. Thus a one-variable polylogarithm
of weight $w$ and depth $d$ has the stated reduction with $K=w-d$.
\end{theorem}
\begin{proof}
On $z<0$, make the substitution $u=(1-t)^{-1}$ in every integration.
Its pullbacks are
\begin{equation}\label{complement:eq:pullbacks}
 \frac{dt}{t}=-\frac{du}{u}-\frac{du}{1-u},\qquad
 \frac{dt}{1-t}=\frac{du}{u}.
\end{equation}
The initial endpoint $t=0$ becomes $u=1$. Multilinearity gives
$H_w(z)=(-1)^K\sum_A J_{\tau_A(w)}(q)$.
Lemma~\ref{complement:lem:connection} proves \eqref{complement:eq:master}.

It remains to check the claimed normal form. The identity
$H_vH_0=H_{v\sh0}$ removes a trailing zero. If a word ends in exactly
$r$ zeros, its coefficient in the shuffle of its prefix with $0$ is
$r$; all other words in that shuffle have fewer trailing zeros.
Induction therefore expresses each $H_u(q)$ as a rational combination
of $L^jH_{u'}(q)$ with $u'$ empty or ending in $1$.
The operation preserves the number of $1$ letters.

For an endpoint word beginning in $0$, the same normalization preserves
its initial $0$ in every nonzero term with no logarithmic factor.
Terms containing a positive power of $\log1$ vanish. The surviving
nonempty terms begin in $0$ and end in $1$, hence are convergent MZVs.
In each split of $\tau_A(w)$ the total number of $1$ letters in the
two factors is $|A|\le K$. This proves the depth bound and the weight
identity. Finally continue analytically from $z<0$ to $D$.
\end{proof}

\begin{corollary}[An argument-enlarged depth bound]\label{complement:cor:halfb}
For $w\ge2$, every one-variable MPL of weight $w$ can be represented
using nonlogarithmic factors of depth at most $\lfloor w/2\rfloor$,
allowing both the original argument $z$ and $q=(1-z)^{-1}$, MZVs, and
powers of $L$ as coefficients. More precisely, use the original
expression when $d\le w/2$, and Theorem~\ref{complement:thm:complement} otherwise.
The all-one case is $L^w/w!$.
\end{corollary}

This corollary is deliberately not phrased as a theorem about arbitrary
multivariable MPLs, or as a reduction inside an unchanged root-of-unity
basket. At $z=i$ the half-point $q$ is not a root of unity. The theorem
explains an argument-enlarged upper bound; it does not prove the absence
or presence of new cyclotomic periods in a smaller comparison algebra.


For the next family put $F_r(z)=\Li_{2,\ones{r-1}}(z)$, as in the
one-two formulas above. Their reciprocal-complementary form is
\[
 F_r(z)=\sum_{k=1}^{r+1}\frac{(-1)^kL^{r+1-k}}{(r+1-k)!}\Li_k(q)
 +(-1)^r\zeta(r+1)-\frac{L^{r+1}}{(r+1)!}.
\]
\section{A double-polylogarithm formula for an all-weight family}

The next formula addresses the first case beyond the classical ladder.
It reduces the height-one family with initial index $3$ to singles and
doubles at $q$, irrespective of its original depth.
Let
\begin{equation}\label{complement:eq:Ck}
 C_1(q)=0,\qquad C_k(q)=\sum_{a=1}^{k-1}\Li_{a,k-a}(q)\quad(k\ge2).
\end{equation}
All MPLs on the right have the one-variable meaning \eqref{complement:eq:mpl}.

\begin{theorem}[Two-zero height-one identity]\label{complement:thm:twozero}
For every $w\ge3$ and $z\in D$,
\begin{align}
\Li_{3,\ones{w-3}}(z)
={}&\frac{L^w}{w!}
 +\sum_{k=1}^{w}\frac{(-1)^kL^{w-k}}{(w-k)!}
   \bigl((k-2)\Li_k(q)+C_k(q)\bigr)\notag\\
&+(-1)^{w-1}\Bigl((L+\Li_1(q))\zeta(w-1)
 +(w-2)\zeta(w)-\zeta(w-1,1)\Bigr).
\label{complement:eq:twozero}
\end{align}
The right side has reduced depth at most two.
\end{theorem}

\begin{proof}
Write $D_q=q\,d/dq=d/dL$. Direct differentiation of the nested series,
followed by continuation, gives
\[
 D_qC_k=C_{k-1}+\frac{q}{1-q}\Li_{k-1}\quad(k\ge2).
\]
Put $V_k=(k-2)\Li_k+C_k$, with $V_1=-\Li_1$.
Then
\[
 D_qV_k=V_{k-1}+\frac{\Li_{k-1}}{1-q}\quad(k\ge2),
 \qquad D_qV_1=-\frac{q}{1-q}.
\]
Differentiating the finite sum in \eqref{complement:eq:twozero} telescopes and yields
\[
 \frac{F_{w-2}(z)-(-1)^{w-2}\zeta(w-1)}{q-1}.
\]
The derivative of the explicit $(L+\Li_1(q))\zeta(w-1)$ term supplies
the missing constant term, because
$D_q(L+\Li_1(q))=1/(1-q)$.
The derivative of the entire right side is therefore
$F_{w-2}(z)/(q-1)$.
The left side has the same derivative, since
$d\Li_{3,\ones{w-3}}(z)/dz=F_{w-2}(z)/z$ and
$dz/z=dL/(q-1)$.

For the endpoint constant, recall the convergent depth-two sum identity
\begin{equation}\label{complement:eq:eulersum}
 \sum_{a=2}^{w-1}\zeta(a,w-a)=\zeta(w),\qquad w\ge3.
\end{equation}
It follows, for example, by subtracting the shuffle and stuffle
expansions of the once-regularized product $\zeta(1)\zeta(w-1)$:
its common divergent term $\zeta(1,w-1)$ cancels, leaving precisely
\eqref{complement:eq:eulersum}. Only one logarithmic divergence is involved, so
no higher regularization correction enters. Equivalently one may use
finite cutoffs and subtract the common harmonic-logarithmic part
before taking the limit.

The convergent colored stuffle identity, for $0<q<1$, gives
\[
 \Li_{1,w-1}(q)-\Li_1(q)\zeta(w-1)
 =-\Li_{w-1,1}(1,q)-\Li_w(q).
\]
The right side tends to $-\zeta(w-1,1)-\zeta(w)$ by dominated
convergence. Combining this with \eqref{complement:eq:eulersum} yields
\[
 \lim_{q\to1^-}\bigl(C_w(q)-\Li_1(q)\zeta(w-1)\bigr)
 =-\zeta(w-1,1).
\]
All positive powers of $L$ multiplying at most logarithmically singular
terms vanish. The remaining constants in \eqref{complement:eq:twozero} cancel,
so its right side tends to zero as $z\to0^-$. The left side does too.
Equality of derivatives and the endpoint limit proves the result;
continue to $D$.
\end{proof}

\begin{remark}[The sum identity and the proof dependency]
Equation~\eqref{complement:eq:eulersum} is the classical depth-two Euler sum
formula, not a new conjecture assumed here. The same two-zero formula
can also be obtained directly from Theorem~\ref{complement:thm:complement} by
expanding its two zero positions and normalizing trailing zeros.
The supplied code verifies the agreement of these two independently
specified formulas through weight twelve.
\end{remark}

At weight six, \eqref{complement:eq:twozero} is a complete depth-two representation
of $\Li_{3,1,1,1}(i)$ and of $\Li_{3,1,1,1}(\rho)$.
In particular a depth-four index does not by itself identify a new
depth-four period in the enlarged comparison algebra.
The endpoint double is harmless as a named convergent constant;
standard Euler reductions can simplify it further. For example,
\[
 \zeta(4,1)=2\zeta(5)-\zeta(2)\zeta(3),\qquad
 \zeta(5,1)=\frac34\zeta(6)-\frac12\zeta(3)^2.
\]
These last simplifications are not required for the theorem.

